\documentclass[aspectratio=169]{beamer}

\usetheme{Madrid}
\usecolortheme{default}

\usepackage[T1]{fontenc}
\usepackage[utf8]{inputenc}
\usepackage{booktabs}
\usepackage{graphicx}
\usepackage{siunitx}

\title[Electricity System Modelling]{Electricity System Modelling}
\subtitle{CO2 Emissions, Storage and Transmission}
\author{John Emanuelsson}
\institute{Miljo och Matematisk Modellering}
\date{\today}

\sisetup{
  round-mode=places,
  round-precision=2,
  group-separator={,},
}

\newcommand{\plotplaceholder}[1]{%
  \fbox{%
    \begin{minipage}[c][0.62\textheight][c]{0.9\linewidth}
      \centering
      Plot file not generated yet:\\[0.5em]
      \scriptsize\nolinkurl{#1}
    \end{minipage}%
  }%
}

\newcommand{\maybeplot}[2][0.95\linewidth]{%
  \IfFileExists{#2}{%
    \includegraphics[width=#1,height=0.68\textheight,keepaspectratio]{#2}%
  }{%
    \plotplaceholder{#2}%
  }%
}

\newcommand{\twoplots}[2]{%
  \begin{columns}[T,onlytextwidth]
    \begin{column}{0.49\textwidth}
      \centering
      \maybeplot[\linewidth]{#1}
    \end{column}
    \begin{column}{0.49\textwidth}
      \centering
      \maybeplot[\linewidth]{#2}
    \end{column}
  \end{columns}
}

\begin{document}

\begin{frame}
  \titlepage
\end{frame}

\begin{frame}{Presentation Roadmap}
  \begin{enumerate}
    \item Model setup and assumptions
    \item Exercise 1: baseline system without CO2 cap, storage, or trade
    \item Exercise 2: 90\% CO2 reduction and batteries
    \item Exercise 3: transmission between countries
    \item Exercise 4: nuclear power
    \item Model limitations and discussion
  \end{enumerate}
\end{frame}

\begin{frame}{Model Scope}
  \begin{itemize}
    \item Countries: Sweden, Denmark, and Germany.
    \item Time horizon: one full year with hourly load balance.
    \item Objective: minimize annualized investment cost plus variable operating and fuel costs.
    \item Technologies: wind, PV, gas, Swedish hydro, batteries, transmission, and nuclear.
    \item Existing capacity: only Swedish hydro remains in 2050.
  \end{itemize}
\end{frame}

\begin{frame}{Core Constraints}
  \begin{itemize}
    \item Demand must be met in every country and every hour.
    \item Wind and PV generation are limited by hourly capacity factors.
    \item Swedish hydro has fixed turbine capacity and a cyclic reservoir balance.
    \item Batteries shift energy over time with storage and charge/discharge limits.
    \item Transmission shifts energy between countries with losses and capacity limits.
    \item CO2 emissions from gas are capped in Exercises 2--4.
  \end{itemize}
\end{frame}

\begin{frame}{Scenario Overview}
  \begin{table}
    \centering
    \begin{tabular}{lllll}
      \toprule
      Exercise & CO2 cap & Batteries & Transmission & Nuclear \\
      \midrule
      1 & No & No & No & No \\
      2a & Yes & No & No & No \\
      2b & Yes & Yes & No & No \\
      3 & Yes & Yes & Yes & No \\
      4 & Yes & Yes & Yes & Yes \\
      \bottomrule
    \end{tabular}
  \end{table}
\end{frame}

\begin{frame}{Average Capacity Factors}
  \begin{table}
    \centering
    \begin{tabular}{lrr}
      \toprule
      Country & Wind & PV \\
      \midrule
      Sweden & 41.0\% & 10.0\% \\
      Denmark & 32.0\% & 11.4\% \\
      Germany & 28.5\% & 13.3\% \\
      \bottomrule
    \end{tabular}
  \end{table}
\end{frame}

\begin{frame}{Exercise 1: Baseline}
  \begin{itemize}
    \item Cost: 37.24 bn EUR/year.
    \item CO2: 138.77 Mton/year.
    \item Wind power is a major generation source in all three countries.
    \item Gas plays a significant role in all three countries, but is smaller in Sweden.
    \item PV is significant except in Sweden, and small in Denmark.
    \item Hydro is significant in Sweden.
  \end{itemize}
\end{frame}

\begin{frame}{Exercise 1: Country Differences}
  \begin{itemize}
    \item PV differences are explained by average capacity factors: Germany has the highest, Sweden the lowest.
    \item Sweden has good wind and existing hydro, so the solution uses both.
    \item Germany has the largest demand, so it builds the largest capacities.
    \item All countries use gas because there is no CO2 cap.
  \end{itemize}
\end{frame}

\begin{frame}{Exercise 1: Capacity and Production}
  \twoplots
    {plots/ex1_no_cap_no_storage_no_trade_capacity.pdf}
    {plots/ex1_no_cap_no_storage_no_trade_production.pdf}
\end{frame}

\begin{frame}{Exercise 1: Germany, Hours 147--651}
  \centering
  \maybeplot{plots/ex1_no_cap_no_storage_no_trade_de_147_651.pdf}
\end{frame}

\begin{frame}{Exercise 2: CO2 Cap and Batteries}
  \begin{itemize}
    \item Exercise 2a: no feasible solution was found.
    \item With a 90\% CO2 cap, demand cannot be met without batteries or transmission.
    \item The main issue is that wind and solar are variable, so demand cannot be met at some time points.
    \item Exercise 2b cost: 64.08 bn EUR/year, increased from 37.24.
    \item Exercise 2b CO2: 13.88 Mton/year, decreased from 138.77.
  \end{itemize}
\end{frame}

\begin{frame}{Exercise 2b: Mix Changes}
  \begin{itemize}
    \item Battery capacity is very large in all countries.
    \item Germany has especially large battery and PV capacity.
    \item Gas capacity is smaller, but still significant.
    \item Production is mostly renewable: PV, wind, and hydro in Sweden.
    \item Gas production is much lower for SE and DK.
    \item Battery discharge is small for SE and DK, but large for DE.
  \end{itemize}
\end{frame}

\begin{frame}{Exercise 2b: Capacity and Production}
  \twoplots
    {plots/ex2b_90pct_co2_cap_batteries_capacity.pdf}
    {plots/ex2b_90pct_co2_cap_batteries_production.pdf}
\end{frame}

\begin{frame}{Exercise 2b: Germany, Hours 147--651}
  \centering
  \maybeplot{plots/ex2b_90pct_co2_cap_batteries_de_147_651.pdf}
\end{frame}

\begin{frame}{Exercise 3: Transmission}
  \begin{itemize}
    \item Cost: 48.52 bn EUR/year, decreased from 64.08.
    \item CO2: 13.88 Mton/year, same as Exercise 2b.
    \item Largest transmission capacity: SE-DE, 46.33 GW.
    \item Largest transmitted energy flow: SE to DE, 131.60 TWh/year.
    \item Transmission reduces the need for local overbuild and battery capacity.
    \item Sweden exports much electricity because it has strong wind and hydro.
  \end{itemize}
\end{frame}

\begin{frame}{Exercise 3: Capacity and Production}
  \twoplots
    {plots/ex3_90pct_co2_cap_batteries_transmission_capacity.pdf}
    {plots/ex3_90pct_co2_cap_batteries_transmission_production.pdf}
\end{frame}

\begin{frame}{Exercise 3: Germany, Hours 147--651}
  \centering
  \maybeplot{plots/ex3_90pct_co2_cap_batteries_transmission_de_147_651.pdf}
\end{frame}

\begin{frame}{Exercise 3: Transmission Capacity and Energy}
  \centering
  \maybeplot{plots/ex3_90pct_co2_cap_batteries_transmission_transmission.pdf}
\end{frame}

\begin{frame}{Exercise 4: Nuclear}
  \begin{itemize}
    \item Cost: 43.52 bn EUR/year, decreased from 48.52.
    \item CO2: 13.88 Mton/year, same as Exercise 3.
    \item Nuclear emerges mainly in Germany: 42.50 GW.
    \item Denmark gets only 0.16 GW nuclear, and Sweden gets none.
    \item Battery usage decreases further, but is still used for variable solar and wind production.
    \item Germany has very high demand and no hydro, so non-variable nuclear generation reduces battery needs.
    \item The result is sensitive to nuclear cost, discount rate, renewable costs, battery cost, transmission cost, and the CO2 cap.
  \end{itemize}
\end{frame}

\begin{frame}{Exercise 4: Capacity and Production}
  \twoplots
    {plots/ex4_90pct_co2_cap_batteries_transmission_nuclear_capacity.pdf}
    {plots/ex4_90pct_co2_cap_batteries_transmission_nuclear_production.pdf}
\end{frame}

\begin{frame}{Exercise 4: Germany, Hours 147--651}
  \centering
  \maybeplot{plots/ex4_90pct_co2_cap_batteries_transmission_nuclear_de_147_651.pdf}
\end{frame}

\begin{frame}{Exercise 4: Transmission Capacity and Energy}
  \centering
  \maybeplot{plots/ex4_90pct_co2_cap_batteries_transmission_nuclear_transmission.pdf}
\end{frame}

\begin{frame}{Model Limitations}
  \begin{itemize}
    \item Stylized electricity-only model with simplified technology representation.
    \item Battery power and energy capacity are tied together as 1 MW = 1 MWh.
    \item Transmission grid is represented by aggregate country-to-country links.
  \end{itemize}
\end{frame}

\begin{frame}{Discussion}
  \begin{itemize}
    \item What drives the main cost differences between scenarios?
    \item Which constraints are most important for the final system configuration?
    \item Which assumptions would most affect the results if changed?
  \end{itemize}
\end{frame}

\end{document}
